\uuid{bQNf}
\titre{Pgcd}
\auteur{bodin ; croizat ; sacre}
\organisation{exo7}
\licence{CC BY-NC-SA 4.0}
\datecreate{2022-03-08}
\niveau{L1}
\module{Algèbre}
\chapitre{Arithmétique dans Z}
\sousChapitre{Pgcd, ppcm, algorithme d'Euclide}
\difficulte{2}
\theme{pgcd}
\exo7qcmid{hilisit.arithmetique.006}

\contenu{
\begin{question}
\qclassification{103.03}
    Pour $a=220$ et $b=60$, quelles sont les affirmations vraies ?
    \begin{answers}
        \bad{$\ppcm(a,b) = 440$.}
        \bad{$440$ est un multiple commun à $a$ et $b$.}
        \good{$10$ est un diviseur commun à $a$ et $b$.}
        \good{$\pgcd(a,b) = 20$.}
    \end{answers}
    \begin{explanations}
        Le plus grand diviseur commun à $a=220$ et $b=60$ est $\pgcd(220,60) = 20$ (on peut l'obtenir via l'algorithme d'Euclide, on en dressant les listes exhaustives des diviseurs communs à $220$ et $60$). Puisque $10$ est un diviseur de $20$, $10$ est bien un diviseur commun à $a$ et $b$ (ce qui se voit sur l'écriture des deux nombres : ils finissent par $0$).\\
    En revanche $440$ n'est pas un multiple de $b=60$ (on a $60 \times 7 = 420$ et $60 \times 8 = 480$). On peut d'ailleurs calculer que $\ppcm(a,b) = \frac{a \times b}{\pgcd(a,b)} = 660$.
    \end{explanations}
\end{question}
}
